Absolutely. I’ll extend it into a **publication-style technical report**, with a stronger scientific structure, explicit algorithm blocks, a source-to-method traceability table, architecture diagrams in TikZ, experiment-flow diagrams, refactoring architecture, reproducibility protocol, and a detailed appendix.

 I also checked the project's associated public material: the project describes HypatiaX as a hybrid LLM + symbolic-regression approach for scientific equation discovery, and its tutorial material describes a 131-equation benchmark and statistical-analysis workflow.  SiMLeng Portfolio+1

 The repository itself should be treated as the primary implementation source; I have deliberately **not invented line numbers** where the repository contents available through search did not expose stable line-numbered source. The extended report therefore identifies files/classes/functions where appropriate and labels implementation-derived observations separately from methodological interpretation.

 LLM-HypatiaX project repository

 \\documentclass\[11pt,a4paper\]{article}

 % ============================================================\
 % PACKAGES\
 % ============================================================

 \\usepackage\[utf8\]{inputenc}\
 \\usepackage\[T1\]{fontenc}\
 \\usepackage{lmodern}\
 \\usepackage\[margin=2.35cm\]{geometry}

 \\usepackage{amsmath}\
 \\usepackage{amssymb}\
 \\usepackage{mathtools}\
 \\usepackage{bm}

 \\usepackage{booktabs}\
 \\usepackage{longtable}\
 \\usepackage{array}\
 \\usepackage{multirow}\
 \\usepackage{tabularx}

 \\usepackage{enumitem}\
 \\usepackage{xcolor}\
 \\usepackage{listings}

 \\usepackage{graphicx}\
 \\usepackage{float}\
 \\usepackage{caption}\
 \\usepackage{subcaption}

 \\usepackage{tikz}\
 \\usetikzlibrary{\
 arrows.meta,\
 positioning,\
 shapes.geometric,\
 shapes.multipart,\
 fit,\
 calc,\
 backgrounds,\
 decorations.pathreplacing\
 }

 \\usepackage{algorithm}\
 \\usepackage{algpseudocode}

 \\usepackage{fancyhdr}\
 \\usepackage{titlesec}\
 \\usepackage{microtype}

 \\usepackage{hyperref}\
 \\usepackage{cleveref}

 % ============================================================\
 % COLOURS\
 % ============================================================

 \\definecolor{hypblue}{RGB}{25,75,125}\
 \\definecolor{hypcyan}{RGB}{35,125,155}\
 \\definecolor{hypgreen}{RGB}{50,120,75}\
 \\definecolor{hyporange}{RGB}{190,110,25}\
 \\definecolor{hypred}{RGB}{165,55,55}\
 \\definecolor{hypgray}{RGB}{90,90,90}\
 \\definecolor{codebg}{RGB}{246,247,249}\
 \\definecolor{boxbg}{RGB}{240,245,250}

 % ============================================================\
 % HYPERREF\
 % ============================================================

 \\hypersetup{\
 colorlinks=true,\
 linkcolor=hypblue,\
 citecolor=hypblue,\
 urlcolor=hypblue,\
 pdftitle={HypatiaX Extended Source-Code and Methodological Analysis},\
 pdfauthor={HypatiaX Code Analysis}\
 }

 % ============================================================\
 % CODE\
 % ============================================================

 \\lstset{\
 backgroundcolor=\\color{codebg},\
 basicstyle=\\ttfamily\\small,\
 keywordstyle=\\color{hypblue}\\bfseries,\
 commentstyle=\\color{hypgreen},\
 stringstyle=\\color{hypred},\
 breaklines=true,\
 breakatwhitespace=false,\
 columns=fullflexible,\
 frame=single,\
 rulecolor=\\color{gray!35},\
 showstringspaces=false,\
 numbers=left,\
 numberstyle=\\tiny\\color{gray},\
 numbersep=7pt\
 }

 % ============================================================\
 % PAGE STYLE\
 % ============================================================

 \\pagestyle{fancy}\
 \\fancyhf{}\
 \\fancyhead\[L\]{HypatiaX Extended Analysis}\
 \\fancyhead\[R\]{Source-Code and Methods}\
 \\fancyfoot\[C\]{\\thepage}

 % ============================================================\
 % CUSTOM COMMANDS\
 % ============================================================

 \\newcommand{\\Rtwo}{R^2}\
 \\newcommand{\\LLM}{\\mathrm{LLM}}\
 \\newcommand{\\NN}{\\mathrm{NN}}\
 \\newcommand{\\OOD}{\\mathrm{OOD}}\
 \\newcommand{\\Hyp}{\\mathrm{HypatiaX}}

 \\newcommand{\\code}\[1\]{\\texttt{#1}}

 % ============================================================\
 % TITLE\
 % ============================================================

 \\title{\
 \\vspace{-1cm}\
 \\textbf{HypatiaX: Extended Source-Code and Methodological Analysis}\\\[0.35cm\]\
 \\large\
 LLM--Symbolic--Neural Formula Discovery,\\\
 Experimental Methodology, Reproducibility and Refactoring\
 }

 \\author{\
 Technical analysis of\
 \\texttt{sednabcn/LLM-HypatiaX-REPRO}\\\
 \\texttt{hypatiax/core} and \\texttt{hypatiax/experiments}\
 }

 \\date{September 2026}

 % ============================================================\
 % DOCUMENT\
 % ============================================================

 \\begin{document}

 \\maketitle

 \\begin{abstract}

 This report presents an extended source-code and methodological analysis of the\
 LLM-HypatiaX reproducibility implementation. The primary scope is the\
 \\texttt{hypatiax/core} and \\texttt{hypatiax/experiments} portions of the\
 repository, with related symbolic utilities considered where they are required\
 to understand the experimental execution path.

 The analysis reconstructs the implemented methodology from source rather than\
 treating HypatiaX as a single abstract algorithm. The codebase contains several\
 related systems: a Pure-LLM symbolic baseline, a DeFi-specific LLM--neural\
 hybrid, a more general scientific/engineering hybrid, an uncertainty-weighted\
 ensemble implementation, adaptive neural training and a collection of\
 benchmark and robustness experiments.

 The common methodological pattern is represented as

 $$
D
\rightarrow
f_{\mathrm{LLM}}
\rightarrow
f_{\mathrm{cal}}
\rightarrow
f_{\mathrm{NN}}
\rightarrow
H(f_{\mathrm{cal}},f_{\mathrm{NN}}),
$$

 where $D$ is a natural-language problem description,\
 $f\_{\\mathrm{LLM}}$ is a symbolic hypothesis produced by an LLM,\
 $f\_{\\mathrm{cal}}$ is a numerically calibrated symbolic expression,\
 $f\_{\\mathrm{NN}}$ is a neural approximation, and $H$ is a selection or\
 ensemble policy.

 The report documents the source components, reconstructs representative\
 algorithms, explains the numerical parameter-fitting stage, analyses adaptive\
 neural training, and maps the principal experimental suites to the underlying\
 methodological questions.

 Particular attention is given to scientific reproducibility. The analysis\
 identifies several important issues: the coexistence of multiple hybrid\
 decision policies, possible use of training performance during model selection,\
 dynamic execution of generated code, hard-coded benchmark-specific cases,\
 LLM API nondeterminism, duplicated implementation and coupling between\
 benchmark orchestration and model execution.

 A refactored architecture is proposed based on six independently testable\
 stages: symbolic discovery, expression validation, symbolic calibration,\
 neural training, validation-based hybrid selection and final evaluation.\
 The proposed architecture also introduces strict train/validation/test/OOD\
 separation, typed symbolic expressions, centralised policies, declarative\
 experiment manifests and immutable provenance.

 \\end{abstract}

 \\tableofcontents\
 \\newpage

 % ============================================================\
 \\section{Executive Summary}\
 % ============================================================

 \\subsection{What HypatiaX implements}

 The implementation is best understood as a \\emph{family of hybrid\
 symbolic--neural formula-discovery systems}.

 Its central idea is to exploit an LLM as a source of semantic and mathematical\
 prior information while retaining numerical learning as a fallback or\
 complement.

 The resulting architecture can be summarised as

 $$
\boxed{
\text{Natural language}
\rightarrow
\text{symbolic hypothesis}
\rightarrow
\text{numerical calibration}
\rightarrow
\text{neural approximation}
\rightarrow
\text{hybrid decision}
}
$$

 The symbolic component is valuable because it can produce a human-readable\
 equation. The neural component is valuable because it can approximate\
 relationships that are difficult to express symbolically.

 \\subsection{Principal implementation variants}

 The source contains at least four important methodological variants:

 \\begin{enumerate}\
 \\item Pure LLM symbolic discovery.\
 \\item DeFi LLM--NN hybrid with explicit symbolic parameter fitting.\
 \\item General-domain LLM--NN hybrid.\
 \\item Residual/uncertainty-weighted LLM--NN ensemble.\
 \\end{enumerate}

 The neural baseline and adaptive configuration mechanism form a fifth\
 supporting component.

 These variants should not be presented in a paper as if they were a single\
 unchanging algorithm.

 \\subsection{Main findings of the source analysis}

 The source analysis identifies the following characteristics.

 \\begin{itemize}\
 \\item LLMs are used as symbolic hypothesis generators.\
 \\item Symbolic constants can be fitted numerically after structural\
 generation.\
 \\item Neural networks provide a flexible numerical approximation.\
 \\item Several hybrid decision rules exist.\
 \\item Domain-specific overrides can alter the normal selection process.\
 \\item Neural architecture and transformations can be selected adaptively.\
 \\item Experiments explicitly examine noise, sample complexity,\
 extrapolation and instability.\
 \\item Reproducibility is addressed through deterministic NN configuration,\
 caching and experiment metadata.\
 \\end{itemize}

 \\subsection{Most important refactoring recommendation}

 The most important architectural change is to replace multiple independent\
 hybrid implementations with a common pipeline:

 $$
\boxed{
Discovery
\rightarrow
Validation
\rightarrow
Calibration
\rightarrow
Training
\rightarrow
Selection
\rightarrow
Evaluation
}
$$

 The model-selection step should operate on a validation set, not on the final\
 test set.

 % ============================================================\
 \\section{Scope and Source-Analysis Method}\
 % ============================================================

 \\subsection{Repository scope}

 The analysis targets:

 \\begin{itemize}\
 \\item \\code{hypatiax/core};\
 \\item \\code{hypatiax/experiments};\
 \\item directly required symbolic utilities.\
 \\end{itemize}

 The repository's public project material describes HypatiaX as a framework\
 for formula discovery using generative AI and symbolic regression. Its\
 associated tutorials describe a benchmark workflow and publication-oriented\
 statistical analysis. These project-level descriptions provide context, while\
 the source code remains the primary evidence for the implementation analysis.

 \\subsection{Evidence hierarchy}

 Three types of evidence are distinguished.

 \\begin{enumerate}\
 \\item \\textbf{Implementation fact:} behaviour directly represented by the\
 source code.\
 \\item \\textbf{Methodological interpretation:} a scientific interpretation\
 of what the implementation is doing.\
 \\item \\textbf{Recommendation:} a proposed improvement that is not itself\
 part of the existing implementation.\
 \\end{enumerate}

 This distinction is important because a source-code implementation can contain\
 engineering decisions that differ from the idealised algorithm described in a\
 paper.

 \\subsection{Static-analysis limitation}

 This report is a static source/methodology analysis. It does not claim that\
 every benchmark has been executed from a clean environment or that every\
 reported numerical result has independently been reproduced.

 Consequently:

 $$
\text{source inspection}
\neq
\text{independent experimental replication}.
$$

 This distinction should remain explicit in any publication or technical\
 appendix derived from this report.

 % ============================================================\
 \\section{Source Inventory}\
 % ============================================================

 \\subsection{Core components}

 \\begin{longtable}{p{0.36\\linewidth}p{0.54\\linewidth}}\
 \\caption{Core implementation inventory}\
 \\label{tab:coreinventory}\\\
 \\toprule\
 \\textbf{Component} & \\textbf{Role} \\\
 \\midrule\
 \\endfirsthead

 \\toprule\
 \\textbf{Component} & \\textbf{Role} \\\
 \\midrule\
 \\endhead

 \\code{base\_pure\_llm/}\
 &\
 Pure LLM formula-discovery baseline.\
 \\

 \\code{generation/hybrid\_defi\_system/}\
 &\
 DeFi-oriented LLM--NN hybrid.\
 \\

 \\code{generation/hybrid\_all\_domains\_llm\_nn/}\
 &\
 General-domain LLM--NN hybrid.\
 \\

 \\code{generation/hybrid\_defi\_llm\_nn/}\
 &\
 Explicit ensemble implementation.\
 \\

 \\code{training/adaptive\_config.py}\
 &\
 A-priori neural configuration.\
 \\

 \\code{training/baseline\_neural\_network\_defi\_improved.py}\
 &\
 Neural baseline/training implementation.\
 \\

 \\bottomrule\
 \\end{longtable}

 \\subsection{Principal experiment components}

 \\begin{longtable}{p{0.40\\linewidth}p{0.50\\linewidth}}\
 \\caption{Experiment inventory}\
 \\label{tab:expinventory}\\\
 \\toprule\
 \\textbf{Experiment} & \\textbf{Purpose} \\\
 \\midrule\
 \\endfirsthead

 \\toprule\
 \\textbf{Experiment} & \\textbf{Purpose} \\\
 \\midrule\
 \\endhead

 \\code{exp1\_ablation.py}\
 &\
 Core ablation and instability/extrapolation analysis.\
 \\

 \\code{run\_comparative\_suite\_benchmark\_v2.py}\
 &\
 Comparative benchmark execution.\
 \\

 \\code{run\_comparative\_suite\_benchmark\_pca.py}\
 &\
 PCA-related comparative experiment.\
 \\

 \\code{run\_instability\_suite.py}\
 &\
 Repeated-run instability analysis.\
 \\

 \\code{run\_noise\_sweep\_benchmark.py}\
 &\
 Noise robustness.\
 \\

 \\code{run\_sample\_complexity\_benchmark.py}\
 &\
 Data-efficiency analysis.\
 \\

 \\code{portfolio\_variance\_v3c2.py}\
 &\
 Portfolio/multi-run variance analysis.\
 \\

 \\code{test\_enhanced\_defi\_extrapolation.py}\
 &\
 Extrapolation regression testing.\
 \\

 \\bottomrule\
 \\end{longtable}

 % ============================================================\
 \\section{System Architecture}\
 % ============================================================

 \\subsection{High-level architecture}

 The current methodology can be represented by the following conceptual\
 architecture.

 \\begin{figure}\[H\]\
 \\centering\
 \\begin{tikzpicture}\[\
 node distance=0.9cm and 0.65cm,\
 box/.style={\
 rectangle,\
 rounded corners,\
 draw=hypblue,\
 fill=boxbg,\
 minimum width=2.7cm,\
 minimum height=1cm,\
 align=center\
 },\
 arrow/.style={\
 -{Latex\[length=2.5mm\]},\
 thick,\
 draw=hypblue\
 }\
 \]

 \\node\[box\] (description) {Scientific /\\financial description};

 \\node\[box, right=of description\] (llm) {LLM\\symbolic discovery};

 \\node\[box, right=of llm\] (symbolic) {Symbolic\\expression};

 \\node\[box, right=of symbolic\] (fit) {Numerical\\parameter fitting};

 \\node\[box, below=of fit\] (nn) {Adaptive\\neural model};

 \\node\[box, left=of nn\] (selection) {Hybrid\\selection};

 \\node\[box, left=of selection\] (evaluation) {Evaluation\\and OOD};

 \\draw\[arrow\] (description) -- (llm);\
 \\draw\[arrow\] (llm) -- (symbolic);\
 \\draw\[arrow\] (symbolic) -- (fit);\
 \\draw\[arrow\] (fit) -- (selection);\
 \\draw\[arrow\] (nn) -- (selection);\
 \\draw\[arrow\] (selection) -- (evaluation);

 \\draw\[arrow\] (description.south) |- (nn.west);

 \\end{tikzpicture}\
 \\caption{Conceptual architecture reconstructed from the implementation.}\
 \\label{fig:architecture}\
 \\end{figure}

 \\subsection{Common abstraction}

 The method can be represented mathematically by:

 $$
\mathcal{D}
=
(D,X,y,M),
$$

 where:

 \\begin{itemize}\
 \\item $D$ is natural-language description;\
 \\item $X$ is the input matrix;\
 \\item $y$ is the target;\
 \\item $M$ is metadata.\
 \\end{itemize}

 The LLM produces:

 $$
f_{\LLM}(X;\theta).
$$

 Parameter fitting produces:

 $$
\hat{\theta}
=
\arg\min_{\theta}
L(y,f_{\LLM}(X;\theta)).
$$

 The symbolic predictor becomes:

 $$
f_S(X)=f_{\LLM}(X;\hat{\theta}).
$$

 The neural model produces:

 $$
f_N(X;\phi).
$$

 The final hybrid system produces:

 $$
f_H(X)
=
H(f_S,f_N).
$$

 % ============================================================\
 \\section{Method A: Pure LLM Symbolic Discovery}\
 % ============================================================

 \\subsection{Source}

 The primary implementation is:

 \\begin{center}\
 \\code{hypatiax/core/base\_pure\_llm/}\\\
 \\code{baseline\_pure\_llm\_defi\_discovery.py}.\
 \\end{center}

 \\subsection{Purpose}

 The Pure-LLM baseline tests whether a language model can infer a mathematical\
 relationship directly from a scientific or financial description.

 The output is intended to be more than a numerical prediction. It may include:

 \\begin{itemize}\
 \\item a formula;\
 \\item LaTeX;\
 \\item executable representation;\
 \\item assumptions;\
 \\item explanatory text;\
 \\item variable interpretation.\
 \\end{itemize}

 \\subsection{Reconstructed algorithm}

 \\begin{algorithm}\[H\]\
 \\caption{Pure-LLM symbolic discovery}\
 \\label{alg:purellm}\
 \\begin{algorithmic}\[1\]

 \\Require Description $D$, variables $V$, metadata $M$, observations $(X,y)$

 \\If{known special case exists}\
 \\State Apply analytical/statistical correction\
 \\EndIf

 \\State Construct prompt $P(D,V,M)$\
 \\State Query LLM using $P$\
 \\State Parse returned response\
 \\State Extract symbolic expression $f$\
 \\State Extract executable representation\
 \\State Evaluate $f(X)$\
 \\State Compute $R^2$, RMSE and MAE\
 \\State \\Return symbolic expression, explanation and metrics

 \\end{algorithmic}\
 \\end{algorithm}

 \\subsection{Analytical corrections}

 The implementation does not rely exclusively on the LLM.

 For selected scientific relationships, numerical or analytical corrections are\
 applied.

 For example, a Michaelis--Menten relationship is:

 $$
v(S)=\frac{V_{\max}S}{K_m+S}.
$$

 An appropriate transformation can allow numerical estimation of $V\_{\\max}$\
 and $K\_m$.

 The broader methodological principle is:

 $$
\text{LLM structural prior}
+
\text{classical numerical estimation}.
$$

 \\subsection{Advantages}

 \\begin{itemize}\
 \\item Directly produces interpretable mathematical expressions.\
 \\item Can use semantic information that numerical-only methods do not see.\
 \\item Can exploit known scientific vocabulary and equation families.\
 \\item Can provide explanatory assumptions alongside predictions.\
 \\end{itemize}

 \\subsection{Limitations}

 \\begin{itemize}\
 \\item Generated expressions can be mathematically plausible but numerically\
 incorrect.\
 \\item LLM responses can vary across API calls.\
 \\item Generated executable code introduces a software/security boundary.\
 \\item Benchmark-specific corrections can reduce the distinction between a\
 general method and a collection of known cases.\
 \\end{itemize}

 % ============================================================\
 \\section{Method B: DeFi LLM--NN Hybrid}\
 % ============================================================

 \\subsection{Source}

 \\begin{center}\
 \\code{hypatiax/core/generation/hybrid\_defi\_system/}\\\
 \\code{hybrid\_system\_nn\_defi\_domain.py}\
 \\end{center}

 \\subsection{Pipeline}

 The implementation contains four principal stages:

 $$
\boxed{
\text{LLM}
\rightarrow
\text{parameter fitting}
\rightarrow
\text{NN}
\rightarrow
\text{hybrid decision}
}
$$

 \\subsection{Stage 1: symbolic generation}

 The LLM proposes:

 $$
f_S(x;\theta_0).
$$

 The formula is evaluated against observations.

 The main predictive metric is:

 $$
R^2
=
1-
\frac{\sum_i(y_i-\hat y_i)^2}
{\sum_i(y_i-\bar y)^2}.
$$

 \\subsection{Stage 2: parameter extraction}

 A key implementation feature is numerical extraction from symbolic formulae.

 Suppose the generated expression is:

 $$
f(x)=\frac{a x}{b+x}.
$$

 The constants can be interpreted as unknown parameters:

 $$
f(x;\theta)
=
\frac{\theta_0x}{\theta_1+x}.
$$

 The fitting problem becomes:

 $$
\hat{\theta}
=
\arg\min_\theta
\sum_i
\left(
y_i-f(x_i;\theta)
\right)^2.
$$

 \\subsection{Stage 3: neural model}

 The neural network uses an MLP architecture with adaptive widths.

 A representative architecture is:

 $$
\mathrm{Linear}(d,h_1)
\rightarrow
\mathrm{LayerNorm}
\rightarrow
\mathrm{SiLU}
\rightarrow
\mathrm{Linear}(h_1,h_2)
\rightarrow
\cdots
\rightarrow
\mathrm{Linear}(h_k,1).
$$

 The number of hidden units depends on the number of variables.

 \\subsection{Stage 4: decision}

 Let:

 $$
R_S=R^2_S
$$

 and:

 $$
R_N=R^2_N.
$$

 The DeFi implementation contains threshold-based logic that can be abstracted\
 as:

 $$
\begin{cases}
S & R_S\text{ is strong and sufficiently exceeds }R_N,\\
H & R_S,R_N\text{ are both adequate and sufficiently close},\\
N & \text{otherwise}.
\end{cases}
$$

 When an ensemble is selected:

 $$
\hat y_H
=
w_S\hat y_S+w_N\hat y_N.
$$

 A score-proportional weighting is approximately:

 $$
w_S=
\frac{R_S}{R_S+R_N},
\qquad
w_N=
\frac{R_N}{R_S+R_N}.
$$

 % ============================================================\
 \\section{Method C: General All-Domain Hybrid}\
 % ============================================================

 \\subsection{Source}

 \\begin{center}\
 \\code{hypatiax/core/generation/hybrid\_all\_domains\_llm\_nn/}\\\
 \\code{hybrid\_system\_llm\_nn\_all\_domains.py}\
 \\end{center}

 \\subsection{Method}

 The all-domain version follows the same general architecture but has a\
 different selection policy.

 The important distinction is:

 $$
\boxed{
\text{same components}
\neq
\text{same hybrid policy}.
}
$$

 The implementation can prefer the symbolic solution when it performs at least\
 as well as the NN, use an ensemble when the NN is sufficiently strong, and use\
 the NN as a fallback when symbolic generation fails.

 \\subsection{Domain override}

 A \\code{force\_llm} mechanism can explicitly favour the symbolic solution for\
 selected domains.

 This can be expressed as:

 $$
\delta_{\mathrm{domain}}=1
\Rightarrow
f_H=f_S.
$$

 This is a legitimate domain prior, but it should be treated as part of the\
 experimental configuration rather than hidden inside a supposedly universal\
 algorithm.

 % ============================================================\
 \\section{Method D: Uncertainty-Weighted Ensemble}\
 % ============================================================

 \\subsection{Source}

 \\begin{center}\
 \\code{hypatiax/core/generation/hybrid\_defi\_llm\_nn/}\\\
 \\code{hybrid\_ensemble\_system\_defi\_domain.py}\
 \\end{center}

 \\subsection{Residual weighting}

 Define residuals:

 $$
e_S=y-\hat y_S,
$$

 $$
e_N=y-\hat y_N.
$$

 Estimate:

 $$
\sigma_S=\operatorname{std}(e_S),
$$

 $$
\sigma_N=\operatorname{std}(e_N).
$$

 Then define:

 $$
q_S=
\frac{\max(R_S,0)}
{\sigma_S+\epsilon},
$$

 $$
q_N=
\frac{\max(R_N,0)}
{\sigma_N+\epsilon}.
$$

 Normalise:

 $$
w_S=\frac{q_S}{q_S+q_N},
\qquad
w_N=\frac{q_N}{q_S+q_N}.
$$

 The final predictor is:

 $$
\hat y_H
=
w_S\hat y_S+w_N\hat y_N.
$$

 \\subsection{Methodological caution}

 This weighting should not automatically be called a calibrated uncertainty\
 estimate.

 If $\\sigma\_S$ and $\\sigma\_N$ are measured on training observations, then:

 $$
\sigma_{\mathrm{train}}
\neq
\sigma_{\mathrm{OOD}}.
$$

 A refactored implementation should estimate ensemble reliability on a\
 validation set.

 % ============================================================\
 \\section{Adaptive Neural Configuration}\
 % ============================================================

 \\subsection{Source}

 \\code{hypatiax/core/training/adaptive\_config.py}

 \\subsection{Motivation}

 The implementation attempts to avoid exhaustive hyperparameter searches.

 Instead:

 $$
C=f(X,y,M,B).
$$

 Here:

 \\begin{itemize}\
 \\item $X$ is the input data;\
 \\item $y$ is the target;\
 \\item $M$ is metadata;\
 \\item $B$ is the available computational budget;\
 \\item $C$ is the training configuration.\
 \\end{itemize}

 \\subsection{Observed signals}

 The configuration can depend on:

 \\begin{itemize}\
 \\item input dimensionality;\
 \\item sample count;\
 \\item target sign;\
 \\item dynamic range;\
 \\item orders of magnitude;\
 \\item pole-like distributions;\
 \\item positivity of inputs;\
 \\item problem difficulty;\
 \\item extrapolation status.\
 \\end{itemize}

 \\subsection{Target transformations}

 For suitable targets:

 $$
y'=\log|y|.
$$

 For suitable positive variables:

 $$
x'_j=\log(x_j).
$$

 The important methodological point is that transformations are selected\
 from observed characteristics rather than through unconstrained tuning.

 \\subsection{Reproducibility}

 The system can derive a seed from the problem description:

 $$
s=
\operatorname{SHA256}(D)\bmod2^{31}.
$$

 Thus:

 $$
D_1=D_2
\Rightarrow
s_1=s_2.
$$

 This provides deterministic neural initialisation for a fixed description.

 % ============================================================\
 \\section{Neural Training Algorithm}\
 % ============================================================

 \\begin{algorithm}\[H\]\
 \\caption{Adaptive neural training}\
 \\label{alg:nn}\
 \\begin{algorithmic}\[1\]

 \\Require Training data $(X,y)$ and metadata $M$

 \\State Extract data characteristics\
 \\State Determine input/output transformations\
 \\State Select architecture based on dimensionality\
 \\State Construct MLP\
 \\State Initialise deterministic seed\
 \\State Scale inputs and target\
 \\State Initialise AdamW optimiser\
 \\State Initialise learning-rate scheduler

 \\For{$t=1,\\ldots,T$}\
 \\State Compute $\\hat y=f\_\\phi(X)$\
 \\State Compute MSE loss\
 \\State Backpropagate\
 \\State Clip gradients\
 \\State Update parameters\
 \\State Update scheduler

```
\If{validation metric improves}
    \State Store best model state
\EndIf

\If{early stopping criterion satisfied}
    \State break
\EndIf
```

 \\EndFor

 \\State Restore best model\
 \\State Return trained network

 \\end{algorithmic}\
 \\end{algorithm}

 % ============================================================\
 \\section{Experimental Framework}\
 % ============================================================

 \\subsection{Experimental questions}

 The experiment directory addresses several different scientific questions.

 \\begin{longtable}{p{0.29\\linewidth}p{0.26\\linewidth}p{0.34\\linewidth}}\
 \\caption{Research questions addressed by the experiment suite}\
 \\label{tab:questions}\\\
 \\toprule\
 \\textbf{Question} & \\textbf{Variable} & \\textbf{Relevant experiment} \\\
 \\midrule\
 \\endfirsthead

 \\toprule\
 \\textbf{Question} & \\textbf{Variable} & \\textbf{Relevant experiment} \\\
 \\midrule\
 \\endhead

 Can the hybrid recover equations?\
 &\
 Method\
 &\
 Core ablation\
 \\

 How does noise affect discovery?\
 &\
 $\\sigma$\
 &\
 Noise sweep\
 \\

 How much data is required?\
 &\
 $N$\
 &\
 Sample-complexity suite\
 \\

 Does the method extrapolate?\
 &\
 OOD range\
 &\
 Extrapolation experiments\
 \\

 How stable is the method?\
 &\
 Random seed/run\
 &\
 Instability suite\
 \\

 How does it compare with other methods?\
 &\
 Method\
 &\
 Comparative suite\
 \\

 \\bottomrule\
 \\end{longtable}

 % ------------------------------------------------------------\
 \\subsection{Core-15 / ablation methodology}

 The ablation suite contains scientific and DeFi equations.

 The experimental structure can be represented as:

 $$
\text{Equation}
\rightarrow
\text{synthetic data}
\rightarrow
\text{noise}
\rightarrow
\text{train/test}
\rightarrow
\text{method execution}
\rightarrow
\text{metrics}.
$$

 Representative domains include:

 \\begin{itemize}\
 \\item Chemistry;\
 \\item Biology;\
 \\item Physics;\
 \\item DeFi automated-market-maker relationships;\
 \\item financial risk.\
 \\end{itemize}

 \\subsection{Noise model}

 A representative perturbation is:

 $$
y_{\mathrm{obs}}
=
y+\epsilon,
$$

 where:

 $$
\epsilon
\sim
\mathcal N
\left(
0,\sigma^2
\right).
$$

 The implementation scales the noise relative to the target distribution in\
 relevant experiments.

 % ------------------------------------------------------------\
 \\subsection{Extrapolation methodology}

 Interpolation evaluates:

 $$
x\in\mathcal X_{\mathrm{train}}.
$$

 Extrapolation evaluates:

 $$
x\notin\mathcal X_{\mathrm{train}}.
$$

 The distinction is fundamental for equation discovery because the objective is\
 often recovery of the underlying law rather than merely approximation within\
 the observed support.

 A symbolic model such as:

 $$
f(x)=ax^2+b
$$

 has a prescribed behaviour outside the observed range.

 A generic MLP:

 $$
f_\phi(x)
$$

 does not necessarily preserve such functional structure outside the training\
 distribution.

 % ------------------------------------------------------------\
 \\subsection{Sample complexity}

 The sample-complexity suite estimates:

 $$
R^2(N)
$$

 for different values of:

 $$
N=|\mathcal X_{\mathrm{train}}|.
$$

 A useful analysis should report:

 $$
\frac{\partial R^2}{\partial N}
$$

 qualitatively through learning curves rather than only comparing a single\
 sample size.

 % ------------------------------------------------------------\
 \\subsection{Instability}

 For repeated runs:

 $$
\mu=
\frac{1}{K}\sum_{k=1}^K R_k^2,
$$

 and:

 $$
\sigma_R=
\sqrt{
\frac{1}{K-1}
\sum_{k=1}^K
(R_k^2-\mu)^2
}.
$$

 The coefficient of variation can additionally be useful:

 $$
CV=
\frac{\sigma_R}{|\mu|+\epsilon}.
$$

 The probability of successful discovery is:

 $$
P_{\tau}
=
\frac{1}{K}
\sum_{k=1}^K
\mathbf{1}[R_k^2>\tau].
$$

 This is more informative than reporting only the best run.

 % ============================================================\
 \\section{Comparative Evaluation}\
 % ============================================================

 A rigorous comparison should distinguish at least four quantities:

 $$
R^2_{\mathrm{train}},
\qquad
R^2_{\mathrm{val}},
\qquad
R^2_{\mathrm{test}},
\qquad
R^2_{\mathrm{OOD}}.
$$

 They answer different questions.

 \\begin{center}\
 \\begin{tabular}{ll}\
 \\toprule\
 Metric & Meaning \\\
 \\midrule\
 $R^2\_{\\mathrm{train}}$ & Fit quality \\\
 $R^2\_{\\mathrm{val}}$ & Selection/generalisation during development \\\
 $R^2\_{\\mathrm{test}}$ & Final interpolation performance \\\
 $R^2\_{\\mathrm{OOD}}$ & Extrapolation/generalisation beyond training support \\\
 \\bottomrule\
 \\end{tabular}\
 \\end{center}

 A model that maximises training $R^2$ should not automatically be considered\
 the best scientific model.

 % ============================================================\
 \\section{Advantages of the HypatiaX Approach}\
 % ============================================================

 \\subsection{Interpretability}

 The most direct advantage is:

 $$
\hat y=\hat f(x_1,\ldots,x_d).
$$

 The result can be inspected as an equation.

 \\subsection{Semantic priors}

 The LLM can exploit natural-language information such as:

 \\begin{itemize}\
 \\item variable names;\
 \\item scientific concepts;\
 \\item known relationships;\
 \\item units;\
 \\item domain terminology.\
 \\end{itemize}

 \\subsection{Structural discovery plus numerical fitting}

 The decomposition:

 $$
\text{structure}
+
\text{parameters}
$$

 is methodologically powerful.

 The LLM can propose the structure:

 $$
f(x;\theta),
$$

 while numerical optimisation estimates:

 $$
\hat{\theta}.
$$

 \\subsection{Flexible numerical fallback}

 The neural model allows the framework to handle cases in which the LLM\
 hypothesis is incomplete.

 \\subsection{Extrapolation-oriented evaluation}

 The explicit OOD experiments are particularly appropriate for symbolic\
 regression.

 \\subsection{Instability analysis}

 Repeated-run analysis recognises that stochastic methods should not be judged\
 solely by their best execution.

 % ============================================================\
 \\section{Limitations and Threats to Validity}\
 % ============================================================

 \\subsection{Internal validity}

 The most important issue is potential leakage between fitting and selection.

 If:

 $$
R_S=R^2(D_{\mathrm{train}})
$$

 and:

 $$
R_N=R^2(D_{\mathrm{train}})
$$

 are used to select the final model, then:

 $$
D_{\mathrm{train}}
$$

 influences both fitting and selection.

 The resulting test may overestimate generalisation if the separation is not\
 strictly maintained.

 \\subsection{Construct validity}

 The LLM is not simply a symbolic-regression engine.

 It is simultaneously:

 \\begin{itemize}\
 \\item a language model;\
 \\item a mathematical prior;\
 \\item a source of candidate functional structures;\
 \\item potentially a source of benchmark-specific prior knowledge.\
 \\end{itemize}

 Therefore the claim being evaluated should be precise:

 $$
\text{LLM-assisted equation discovery}
$$

 rather than simply:

 $$
\text{symbolic regression}.
$$

 \\subsection{External validity}

 A benchmark dominated by synthetic equations may not represent real-world\
 scientific measurement.

 Real data can contain:

 \\begin{itemize}\
 \\item missing variables;\
 \\item measurement error;\
 \\item correlated noise;\
 \\item hidden confounders;\
 \\item unit inconsistencies;\
 \\item regime changes.\
 \\end{itemize}

 These should be considered in future validation.

 \\subsection{LLM reproducibility}

 Remote model behaviour may change because of:

 \\begin{itemize}\
 \\item model version;\
 \\item provider changes;\
 \\item prompt changes;\
 \\item system-message changes;\
 \\item API configuration;\
 \\item backend nondeterminism.\
 \\end{itemize}

 A frozen-response cache is therefore useful for exact reproduction.

 \\subsection{Code-execution risk}

 A generated Python expression should not be treated as a trusted mathematical\
 expression.

 A safer architecture is:

 $$
\text{LLM}
\rightarrow
\text{JSON/DSL}
\rightarrow
\text{parser}
\rightarrow
\text{validator}
\rightarrow
\text{evaluator}.
$$

 % ============================================================\
 \\section{Refactoring Architecture}\
 % ============================================================

 \\subsection{Target package structure}

 The proposed structure is:

 \\begin{lstlisting}\
 hypatiax/\
 |\
 +-- discovery/\
 | +-- llm.py\
 | +-- prompts.py\
 | +-- parser.py\
 | +-- cache.py\
 |\
 +-- symbolic/\
 | +-- expression.py\
 | +-- validator.py\
 | +-- evaluator.py\
 | +-- parameter\_fit.py\
 |\
 +-- models/\
 | +-- neural.py\
 | +-- trainer.py\
 | +-- adaptive\_config.py\
 |\
 +-- hybrid/\
 | +-- selector.py\
 | +-- ensemble.py\
 | +-- policies.py\
 |\
 +-- evaluation/\
 | +-- metrics.py\
 | +-- splits.py\
 | +-- extrapolation.py\
 | +-- instability.py\
 |\
 +-- experiments/\
 | +-- definitions/\
 | +-- runners/\
 | +-- analyses/\
 | +-- reports/\
 |\
 +-- provenance/\
 | +-- schema.py\
 | +-- recorder.py\
 \\end{lstlisting}

 % ------------------------------------------------------------\
 \\subsection{Typed expression layer}

 The symbolic expression should use a controlled AST.

 For example:

 $$
E::=
x_i
\mid
c
\mid
E+E
\mid
E-E
\mid
E E
\mid
E^E
\mid
\log E
\mid
\exp E.
$$

 The LLM should never need to emit arbitrary Python.

 % ------------------------------------------------------------\
 \\subsection{Canonical model interface}

 All model classes should implement:

 \\begin{lstlisting}\[language=Python\]\
 class FormulaModel:

```
def fit(self, X, y, metadata):
    ...

def predict(self, X):
    ...

def describe(self):
    ...

def provenance(self):
    ...
```

 \\end{lstlisting}

 This allows all methods to be evaluated identically.

 % ------------------------------------------------------------\
 \\subsection{Canonical result object}

 A result should contain:

 \\begin{lstlisting}\
 {\
 "formula": "...",\
 "parameters": {...},\
 "model\_type": "...",\
 "train\_metrics": {...},\
 "validation\_metrics": {...},\
 "test\_metrics": {...},\
 "ood\_metrics": {...},\
 "selection": "...",\
 "provenance": {...}\
 }\
 \\end{lstlisting}

 % ------------------------------------------------------------\
 \\subsection{Central selection policy}

 Instead of embedding thresholds inside each hybrid class:

 \\begin{lstlisting}\[language=Python\]\
 class HybridPolicy:

```
def select(
    self,
    symbolic_result,
    neural_result,
    validation_data
):
    raise NotImplementedError
```

 \\end{lstlisting}

 Concrete implementations become:

 \\begin{lstlisting}\
 ValidationR2Policy\
 MarginPolicy\
 UncertaintyPolicy\
 FixedBlendPolicy\
 SymbolicOnlyPolicy\
 NeuralOnlyPolicy\
 \\end{lstlisting}

 This makes experimental comparison much cleaner.

 % ============================================================\
 \\section{Proposed Refactored Algorithm}\
 % ============================================================

 \\begin{algorithm}\[H\]\
 \\caption{Canonical Refactored HypatiaX}\
 \\label{alg:canonical}\
 \\begin{algorithmic}\[1\]

 \\Require\
 Description $D$,\
 training set $D\_T$,\
 validation set $D\_V$,\
 test set $D\_E$,\
 OOD set $D\_O$,\
 metadata $M$

 \\Statex

 \\State \\textbf{Discovery}\
 \\State $P\\gets\\textsc{BuildPrompt}(D,M)$\
 \\State $r\\gets\\textsc{LLM}(P)$\
 \\State $e\\gets\\textsc{Parse}(r)$\
 \\State $\\textsc{Validate}(e)$

 \\Statex

 \\State \\textbf{Symbolic calibration}\
 \\State $\\hat{\\theta}\\gets\
 \\textsc{FitParameters}(e,D\_T)$\
 \\State $f\_S\\gets\\textsc{Instantiate}(e,\\hat{\\theta})$

 \\Statex

 \\State \\textbf{Neural training}\
 \\State $C\\gets\\textsc{AdaptiveConfig}(D\_T,M)$\
 \\State $f\_N\\gets\\textsc{TrainNN}(D\_T,C)$

 \\Statex

 \\State \\textbf{Validation}\
 \\State $m\_S\\gets\\textsc{Evaluate}(f\_S,D\_V)$\
 \\State $m\_N\\gets\\textsc{Evaluate}(f\_N,D\_V)$

 \\Statex

 \\State \\textbf{Selection}\
 \\State $H\\gets\
 \\textsc{Policy}(m\_S,m\_N,D\_V)$

 \\Statex

 \\State \\textbf{Final evaluation}\
 \\State $m\_E\\gets\\textsc{Evaluate}(H,D\_E)$\
 \\State $m\_O\\gets\\textsc{Evaluate}(H,D\_O)$

 \\Statex

 \\State \\textbf{Provenance}\
 \\State $\\textsc{Store}(D,P,r,e,\\hat{\\theta},C,H,m\_E,m\_O)$

 \\State \\Return $H,m\_E,m\_O,\\text{provenance}$

 \\end{algorithmic}\
 \\end{algorithm}

 % ============================================================\
 \\section{Experiment Refactoring}\
 % ============================================================

 \\subsection{Current problem}

 Large experiment scripts frequently combine:

 $$
\text{configuration}
+
\text{data generation}
+
\text{training}
+
\text{evaluation}
+
\text{statistics}
+
\text{plotting}
+
\text{file output}.
$$

 This violates separation of concerns.

 \\subsection{Recommended manifest}

 A declarative experiment can instead be described by:

 \\begin{lstlisting}\
 experiment:\
 name: core15\_hybrid\_ablation

 dataset:\
 equation\_set: core15\
 samples: 200\
 noise: 0.05

 split:\
 train: 0.70\
 validation: 0.15\
 test: 0.15

 ood:\
 enabled: true\
 multiplier: 5

 methods:\
 \- pure\_llm\
 \- neural\
 \- hybrid\
 \- symbolic

 hybrid:\
 policy: validation\_margin

 evaluation:\
 metrics:\
 \- r2\
 \- rmse\
 \- mae

 repetitions:\
 seeds:\
 \- 42\
 \- 99\
 \- 123\
 \- 777\
 \- 2024\
 \\end{lstlisting}

 A single runner can consume this configuration.

 % ============================================================\
 \\section{Reproducibility Protocol}\
 % ============================================================

 \\subsection{Required artefacts}

 Every experiment should save:

 \\begin{enumerate}\
 \\item Git commit hash.\
 \\item Dataset hash.\
 \\item Configuration hash.\
 \\item Prompt hash.\
 \\item Raw LLM response.\
 \\item Parsed symbolic expression.\
 \\item Fitted parameters.\
 \\item Neural configuration.\
 \\item Random seeds.\
 \\item Environment/dependency information.\
 \\item Final metrics.\
 \\item OOD metrics.\
 \\end{enumerate}

 \\subsection{Provenance schema}

 A minimal record is:

 \\begin{lstlisting}\
 {\
 "experiment\_id": "...",\
 "equation\_id": "...",\
 "git\_commit": "...",\
 "dataset\_hash": "...",\
 "config\_hash": "...",\
 "prompt\_hash": "...",\
 "llm\_model": "...",\
 "llm\_response\_hash": "...",\
 "seed": 42,

```
"symbolic_formula": "...",
"fitted_parameters": {},

"neural_config": {},

"selection_policy": "...",

"metrics": {
    "train": {},
    "validation": {},
    "test": {},
    "ood": {}
}
```

 }\
 \\end{lstlisting}

 % ============================================================\
 \\section{Statistical Reporting}\
 % ============================================================

 For $K$ repetitions, report:

 $$
\bar R^2=
\frac{1}{K}
\sum_{k=1}^{K}
R_k^2.
$$

 Report uncertainty:

 $$
s_R=
\sqrt{
\frac{1}{K-1}
\sum_{k=1}^{K}
(R_k^2-\bar R^2)^2
}.
$$

 A $95%$ confidence interval for the mean can be reported as:

 $$
\bar R^2
\pm
t_{0.975,K-1}
\frac{s_R}{\sqrt K}.
$$

 Success probability:

 $$
p_\tau=
\frac{1}{K}
\sum_{k=1}^{K}
\mathbf{1}(R_k^2>\tau).
$$

 For comparing methods across equations, paired statistical analysis should\
 be preferred because the same equations can be evaluated under multiple\
 methods.

 % ============================================================\
 \\section{Detailed Traceability Matrix}\
 % ============================================================

 \\begin{longtable}{p{0.25\\linewidth}p{0.27\\linewidth}p{0.37\\linewidth}}\
 \\caption{Source-to-method traceability}\
 \\label{tab:traceability}\\\
 \\toprule\
 \\textbf{Methodological stage} &\
 \\textbf{Implementation area} &\
 \\textbf{Observed responsibility}\
 \\\
 \\midrule\
 \\endfirsthead

 \\toprule\
 \\textbf{Methodological stage} &\
 \\textbf{Implementation area} &\
 \\textbf{Observed responsibility}\
 \\\
 \\midrule\
 \\endhead

 Problem interpretation\
 &\
 LLM prompt construction\
 &\
 Translate scientific/financial context into a symbolic discovery request.\
 \\

 Symbolic discovery\
 &\
 Pure LLM and hybrid generation\
 &\
 Request and parse candidate mathematical expressions.\
 \\

 Symbolic execution\
 &\
 Hybrid formula evaluation\
 &\
 Evaluate candidate formula against numerical observations.\
 \\

 Parameter calibration\
 &\
 DeFi hybrid\
 &\
 Optimise numerical constants while preserving symbolic structure.\
 \\

 Neural configuration\
 &\
 \\code{adaptive\_config.py}\
 &\
 Select transformations, architecture and training settings.\
 \\

 Neural training\
 &\
 Neural baseline/hybrid\
 &\
 Fit MLP approximation to numerical observations.\
 \\

 Model selection\
 &\
 Hybrid implementations\
 &\
 Compare symbolic and neural performance and choose/blend.\
 \\

 Robustness\
 &\
 Noise/sample suites\
 &\
 Measure degradation under noise and reduced sample size.\
 \\

 Extrapolation\
 &\
 Extrapolation tests\
 &\
 Evaluate behaviour outside the training support.\
 \\

 Instability\
 &\
 Instability suite\
 &\
 Aggregate repeated-run performance and variance.\
 \\

 Reproducibility\
 &\
 Caches/seeds/configuration\
 &\
 Record or control experimental sources of variation.\
 \\

 \\bottomrule\
 \\end{longtable}

 % ============================================================\
 \\section{Recommended Publication Figures}\
 % ============================================================

 A publication-quality version of the work should contain at least the following\
 figures.

 \\subsection{Figure 1: HypatiaX architecture}

 $$
D
\rightarrow
LLM
\rightarrow
Symbolic
\rightarrow
Calibration
\rightarrow
NN
\rightarrow
Selection
\rightarrow
Evaluation.
$$

 \\subsection{Figure 2: Data separation}

 \\begin{tikzpicture}\[\
 box/.style={\
 rectangle,\
 draw=hypblue,\
 rounded corners,\
 minimum width=3cm,\
 minimum height=1cm,\
 align=center\
 }\
 \]

 \\node\[box,fill=blue!8\] (train) {Training\\fit parameters/models};\
 \\node\[box,fill=green!8,right=0.8cm of train\] (val)\
 {Validation\\select model};\
 \\node\[box,fill=orange!12,right=0.8cm of val\] (test)\
 {Test\\final metric};\
 \\node\[box,fill=red!10,right=0.8cm of test\] (ood)\
 {OOD\\extrapolation};

 \\draw\[-{Latex}\] (train) -- (val);\
 \\draw\[-{Latex}\] (val) -- (test);\
 \\draw\[-{Latex}\] (test) -- (ood);

 \\end{tikzpicture}

 \\subsection{Figure 3: Hybrid selection}

 A useful empirical figure should plot:

 $$
R^2_{\mathrm{symbolic}}
$$

 against:

 $$
R^2_{\mathrm{NN}},
$$

 with points coloured by the selected policy.

 \\subsection{Figure 4: Extrapolation}

 Plot:

 $$
R^2_{\mathrm{test}}
$$

 against:

 $$
R^2_{\mathrm{OOD}}.
$$

 This distinguishes interpolation from law recovery.

 \\subsection{Figure 5: Instability}

 Use boxplots or violin plots of:

 $$
R^2_1,\ldots,R^2_K
$$

 for each method and equation family.

 % ============================================================\
 \\section{Recommended Ablation Matrix}\
 % ============================================================

 The following ablation matrix would provide a cleaner causal analysis of the\
 method.

 \\begin{center}\
 \\begin{tabular}{llll}\
 \\toprule\
 Variant & LLM & Calibration & NN \\\
 \\midrule\
 Pure LLM & Yes & No & No \\\
 LLM + calibration & Yes & Yes & No \\\
 NN only & No & No & Yes \\\
 LLM + NN & Yes & No & Yes \\\
 LLM + calibration + NN & Yes & Yes & Yes \\\
 Hybrid without domain prior & Yes & Yes & Yes \\\
 Hybrid with domain prior & Yes & Yes & Yes \\\
 \\bottomrule\
 \\end{tabular}\
 \\end{center}

 The comparison isolates whether performance comes from:

 \\begin{itemize}\
 \\item the LLM;\
 \\item parameter fitting;\
 \\item neural approximation;\
 \\item their combination;\
 \\item domain-specific prior knowledge.\
 \\end{itemize}

 % ============================================================\
 \\section{Refactoring Roadmap}\
 % ============================================================

 A practical migration can be divided into stages.

 \\subsection{Phase 1: Extract shared utilities}

 Move duplicated implementations of:

 \\begin{itemize}\
 \\item $R^2$;\
 \\item RMSE;\
 \\item MAE;\
 \\item formula evaluation;\
 \\item scaling;\
 \\item seed management.\
 \\end{itemize}

 into shared modules.

 \\subsection{Phase 2: Introduce typed symbolic expressions}

 Replace dynamic source execution with an AST/DSL.

 \\subsection{Phase 3: Separate model fitting and selection}

 Create:

 $$
fit()
$$

 and:

 $$
select()
$$

 as separate operations.

 \\subsection{Phase 4: Introduce validation data}

 Move hybrid decision logic from training metrics to validation metrics.

 \\subsection{Phase 5: Convert experiments to manifests}

 Create a generic runner.

 \\subsection{Phase 6: Add provenance}

 Make every output self-describing and reproducible.

 \\subsection{Phase 7: Delete duplicated hybrid implementations}

 The original DeFi and all-domain systems can become configurations of the same\
 underlying engine.

 % ============================================================\
 \\section{Target Refactored Architecture}\
 % ============================================================

 \\begin{figure}\[H\]\
 \\centering

 \\begin{tikzpicture}\[\
 module/.style={\
 rectangle,\
 rounded corners,\
 draw=hypblue,\
 fill=blue!5,\
 minimum width=3.2cm,\
 minimum height=0.85cm,\
 align=center\
 },\
 process/.style={\
 rectangle,\
 rounded corners,\
 draw=hypgreen,\
 fill=green!6,\
 minimum width=3.2cm,\
 minimum height=0.85cm,\
 align=center\
 },\
 arrow/.style={\
 -{Latex\[length=2.5mm\]},\
 thick\
 }\
 \]

 \\node\[module\] (llm) {LLM Discovery};\
 \\node\[module,right=1cm of llm\] (parser) {Expression Parser};\
 \\node\[module,right=1cm of parser\] (validator) {Expression Validator};

 \\node\[process,below=1.1cm of parser\] (fit)\
 {Symbolic Parameter Fitting};

 \\node\[process,left=1cm of fit\] (nn)\
 {Neural Training};

 \\node\[process,right=1cm of fit\] (selector)\
 {Validation Policy};

 \\node\[module,below=1.1cm of fit\] (eval)\
 {Test / OOD Evaluation};

 \\draw\[arrow\] (llm) -- (parser);\
 \\draw\[arrow\] (parser) -- (validator);\
 \\draw\[arrow\] (validator) |- (fit);\
 \\draw\[arrow\] (fit) -- (selector);\
 \\draw\[arrow\] (nn) -- (selector);\
 \\draw\[arrow\] (selector) |- (eval);

 \\end{tikzpicture}

 \\caption{Proposed refactored architecture.}\
 \\label{fig:refactored}\
 \\end{figure}

 % ============================================================\
 \\section{Scientific Interpretation}\
 % ============================================================

 The central scientific hypothesis behind the architecture is that symbolic\
 priors and flexible numerical approximation have complementary strengths.

 Let:

 $$
f^*(x)
$$

 be the unknown true relationship.

 The symbolic learner searches over a structured family:

 $$
f_S(x;\theta)
\in
\mathcal F_S.
$$

 The neural learner searches over a flexible function family:

 $$
f_N(x;\phi)
\in
\mathcal F_N.
$$

 Typically:

 $$
\mathcal F_S
\subsetneq
\mathcal F_N
$$

 in terms of representational flexibility, but symbolic models impose stronger\
 structural constraints.

 The hybrid therefore attempts to exploit:

 $$
\text{bias}
+
\text{flexibility}.
$$

 The scientific value of the approach is greatest when the symbolic prior is\
 informative but imperfect.

 % ============================================================\
 \\section{Key Methodological Distinctions}\
 % ============================================================

 \\subsection{Prediction versus equation discovery}

 A neural model can minimise:

 $$
\sum_i
(y_i-\hat y_i)^2
$$

 without discovering the correct equation.

 Equation discovery additionally asks whether:

 $$
\hat f
\approx
f^*
$$

 as a functional object.

 Therefore:

 $$
\text{low prediction error}
\not\Rightarrow
\text{correct scientific law}.
$$

 \\subsection{Interpolation versus extrapolation}

 Similarly:

 $$
R^2_{\mathrm{test}}\approx1
$$

 does not imply:

 $$
R^2_{\mathrm{OOD}}\approx1.
$$

 This makes the repository's extrapolation experiments particularly relevant\
 to the scientific interpretation of the method.

 \\subsection{Selection versus ensemble}

 There are two fundamentally different hybrid concepts.

 Selection:

 $$
f_H=
\begin{cases}
f_S & \text{if symbolic is selected},\\
f_N & \text{otherwise}.
\end{cases}
$$

 Ensembling:

 $$
f_H=
wf_S+(1-w)f_N.
$$

 These should be reported as different experimental conditions.

 % ============================================================\
 \\section{Final Assessment}\
 % ============================================================

 The implementation represents a technically interesting integration of:

 $$
\boxed{
LLM
+
symbolic\ regression
+
numerical\ optimisation
+
neural\ approximation
}
$$

 The source code demonstrates that the method is substantially more than an LLM\
 prompt that asks for an equation. It includes numerical parameter fitting,\
 adaptive neural training, model comparison, domain-specific logic and a broad\
 experimental programme.

 At the same time, the current implementation should not be described as one\
 single deterministic algorithm.

 The most important methodological distinctions are:

 \\begin{enumerate}\
 \\item Pure LLM versus hybrid discovery.\
 \\item Raw symbolic formula versus calibrated symbolic formula.\
 \\item Neural approximation versus symbolic model.\
 \\item Selection versus ensemble.\
 \\item Interpolation versus extrapolation.\
 \\item Training performance versus validation performance.\
 \\item Generic method versus benchmark-specific correction.\
 \\end{enumerate}

 The principal refactoring recommendation is therefore to make these dimensions\
 explicit rather than allowing them to remain distributed across multiple\
 source files.

 % ============================================================\
 \\section{Conclusions}\
 % ============================================================

 This extended source-code analysis leads to six principal conclusions.

 \\subsection\*{Conclusion 1: HypatiaX is a method family}

 The repository contains multiple hybrid implementations and decision policies.\
 The exact algorithm should therefore always identify its variant.

 \\subsection\*{Conclusion 2: Symbolic calibration is a central contribution}

 The transition:

 $$
f_{\LLM}(x;\theta_0)
\rightarrow
f_{\LLM}(x;\hat\theta)
$$

 is an important bridge between language-model-generated structure and\
 data-driven numerical estimation.

 \\subsection\*{Conclusion 3: The NN is complementary}

 The neural network is not merely a competing baseline. Within the hybrid\
 architecture it serves as a flexible numerical alternative and, in some\
 implementations, an ensemble component.

 \\subsection\*{Conclusion 4: Extrapolation is essential}

 For equation discovery, OOD evaluation is more informative than interpolation\
 alone.

 \\subsection\*{Conclusion 5: Reproducibility requires provenance}

 LLM responses, prompts, model identifiers, code versions, datasets,\
 configurations and random seeds should all be retained.

 \\subsection\*{Conclusion 6: Refactoring can substantially improve scientific clarity}

 The proposed architecture:

 $$
\boxed{
Discovery
\rightarrow
Parsing
\rightarrow
Validation
\rightarrow
Calibration
\rightarrow
Neural\ Training
\rightarrow
Validation\ Selection
\rightarrow
Test/OOD
}
$$

 would preserve the scientific idea while making the implementation easier to\
 test, reproduce and extend.

 % ============================================================\
 \\appendix\
 % ============================================================

 \\section{Appendix A: Source-to-Algorithm Mapping}

 \\begin{longtable}{p{0.38\\linewidth}p{0.50\\linewidth}}\
 \\caption{Detailed source-to-algorithm mapping}\
 \\\
 \\toprule\
 \\textbf{Source component} & \\textbf{Algorithmic role}\\\
 \\midrule\
 \\endfirsthead

 \\toprule\
 \\textbf{Source component} & \\textbf{Algorithmic role}\\\
 \\midrule\
 \\endhead

 \\code{baseline\_pure\_llm\_defi\_discovery.py}\
 &\
 Generate and evaluate symbolic hypotheses from LLM output.\
 \\

 \\code{hybrid\_system\_nn\_defi\_domain.py}\
 &\
 Combine symbolic generation, parameter fitting, NN training and hybrid\
 selection.\
 \\

 \\code{hybrid\_system\_llm\_nn\_all\_domains.py}\
 &\
 Generalise hybrid logic to broader scientific/engineering domains.\
 \\

 \\code{hybrid\_ensemble\_system\_defi\_domain.py}\
 &\
 Construct weighted symbolic/neural prediction ensembles.\
 \\

 \\code{adaptive\_config.py}\
 &\
 Map dataset characteristics to neural training configuration.\
 \\

 \\code{baseline\_neural\_network\_defi\_improved.py}\
 &\
 Provide adaptive NN training and evaluation.\
 \\

 \\code{exp1\_ablation.py}\
 &\
 Compare methods over controlled equation/data conditions.\
 \\

 \\code{run\_instability\_suite.py}\
 &\
 Aggregate repeated runs into stability statistics.\
 \\

 \\code{run\_noise\_sweep\_benchmark.py}\
 &\
 Measure sensitivity to noise.\
 \\

 \\code{run\_sample\_complexity\_benchmark.py}\
 &\
 Measure sensitivity to sample count.\
 \\

 \\code{test\_enhanced\_defi\_extrapolation.py}\
 &\
 Test extrapolation behaviour.\
 \\

 \\bottomrule\
 \\end{longtable}

 % ============================================================\
 \\section{Appendix B: Recommended Reporting Table}\
 % ============================================================

 A final paper should report results in a table of the following form.

 \\begin{center}\
 \\small\
 \\begin{tabular}{lrrrrrr}\
 \\toprule\
 Equation &\
 Train $R^2$ &\
 Val. $R^2$ &\
 Test $R^2$ &\
 OOD $R^2$ &\
 RMSE &\
 MAE\
 \\\
 \\midrule\
 Eq. 1 & -- & -- & -- & -- & -- & --\\\
 Eq. 2 & -- & -- & -- & -- & -- & --\\\
 Eq. 3 & -- & -- & -- & -- & -- & --\\\
 $\\vdots$ & & & & & &\\\
 \\bottomrule\
 \\end{tabular}\
 \\end{center}

 For repeated runs, each entry should preferably be:

 $$
\mu\pm\sigma.
$$

 % ============================================================\
 \\section{Appendix C: Recommended Ablation Results}\
 % ============================================================

 The final experimental paper should isolate:

 $$
\begin{array}{c}
LLM\\
LLM+\text{calibration}\\
NN\\
LLM+NN\\
LLM+\text{calibration}+NN
\end{array}
$$

 so that the contribution of each component can be identified.

 A particularly important comparison is:

 $$
\Delta_{\mathrm{cal}}
=
R^2_{\mathrm{LLM+cal}}
-
R^2_{\mathrm{LLM}},
$$

 and:

 $$
\Delta_{\mathrm{hybrid}}
=
R^2_{\mathrm{hybrid}}
-
\max
\left(
R^2_{\mathrm{LLM+cal}},
R^2_{\mathrm{NN}}
\right).
$$

 These quantities answer two different questions:

 \\begin{itemize}\
 \\item Does numerical calibration improve the symbolic hypothesis?\
 \\item Does combining symbolic and neural models improve upon the strongest\
 individual component?\
 \\end{itemize}

 % ============================================================\
 \\section{Appendix D: Security Recommendation}\
 % ============================================================

 The use of generated executable Python should be considered a security\
 boundary.

 The preferred mathematical representation is a restricted DSL:

 $$
E::=
x_i
\mid
c
\mid
E+E
\mid
E-E
\mid
E\times E
\mid
E^E
\mid
\log E
\mid
\exp E
\mid
\sin E
\mid
\cos E.
$$

 An expression validator should reject:

 \\begin{itemize}\
 \\item imports;\
 \\item attribute access;\
 \\item function definitions;\
 \\item file operations;\
 \\item subprocess execution;\
 \\item arbitrary Python statements.\
 \\end{itemize}

 The numerical evaluator should operate only on the validated expression tree.

 % ============================================================\
 \\section{Appendix E: Reproducibility Checklist}\
 % ============================================================

 Before publication, verify:

 \\begin{itemize}\
 \\item \[ \] Exact Git commit recorded.\
 \\item \[ \] Python version recorded.\
 \\item \[ \] Dependency versions recorded.\
 \\item \[ \] LLM provider recorded.\
 \\item \[ \] LLM model identifier recorded.\
 \\item \[ \] Prompt stored.\
 \\item \[ \] Raw LLM response stored.\
 \\item \[ \] Dataset generation seed recorded.\
 \\item \[ \] Dataset hash recorded.\
 \\item \[ \] Neural seed recorded.\
 \\item \[ \] Training configuration stored.\
 \\item \[ \] Hybrid selection policy stored.\
 \\item \[ \] Train/validation/test/OOD split stored.\
 \\item \[ \] Final test data never used for selection.\
 \\item \[ \] OOD data never used for fitting.\
 \\item \[ \] All repeated runs stored.\
 \\item \[ \] Failed runs retained rather than silently discarded.\
 \\item \[ \] Mean and standard deviation reported.\
 \\item \[ \] Extrapolation reported separately.\
 \\item \[ \] Benchmark-specific corrections documented.\
 \\end{itemize}

 % ============================================================\
 \\section{Appendix F: Minimal Compilation}\
 % ============================================================

 Save this file as:

 \\begin{center}\
 \\code{hypatiax\_extended\_report.tex}\
 \\end{center}

 Compile twice:

 \\begin{lstlisting}\[language=bash\]\
 pdflatex hypatiax\_extended\_report.tex\
 pdflatex hypatiax\_extended\_report.tex\
 \\end{lstlisting}

 For a cleaner automated build:

 \\begin{lstlisting}\[language=bash\]\
 latexmk -pdf -interaction=nonstopmode \
 hypatiax\_extended\_report.tex\
 \\end{lstlisting}

 % ============================================================\
 \\section{Appendix G: Suggested Bibliography}\
 % ============================================================

 The final publication version should include, at minimum:

 \\begin{itemize}\
 \\item the original HypatiaX/LLM-HypatiaX paper;\
 \\item the repository/software release;\
 \\item the symbolic-regression method used as a baseline;\
 \\item the neural-network framework used by the implementation;\
 \\item the LLM/API model documentation corresponding to the experiments;\
 \\item relevant equation-discovery and scientific machine-learning literature.\
 \\end{itemize}

 For the repository itself, cite the exact commit used for the analysis rather\
 than only citing a moving branch.

 \\end{document}

 ### What is different in this extended version

 This version goes beyond the earlier report by adding:

 - a **publication-style abstract and executive summary**;
- a formal **scope/evidence hierarchy**;
- a **source inventory and traceability matrix**;
- reconstructed **algorithm environments**;
- mathematical formulation of the symbolic calibration;
- separate treatment of the **DeFi hybrid**, **all-domain hybrid**, and\
   **uncertainty ensemble**;
- **TikZ architecture diagrams**;
- train/validation/test/OOD methodology;
- explicit **threats to validity**;
- a proposed **typed symbolic AST/DSL**;
- a complete **refactored architecture**;
- a canonical **refactored HypatiaX algorithm**;
- statistical reporting equations;
- a proposed **ablation matrix**;
- a **refactoring roadmap**;
- a **reproducibility/provenance schema**;
- security analysis of generated-code execution;
- a publication-oriented **reproducibility checklist**;
- appendices for source mapping, experimental reporting and ablations.

 The project's own public tutorial material also indicates that its benchmark workflow is intended to support reproduction of a **131-equation benchmark** and subsequent statistical/publication-figure analysis, which is why the extended report separates the methodological source analysis from claims about experimentally reproduced numerical results.  SiMLeng Portfolio

 **Important:** this is still a source-analysis report, not a claim that I independently executed all 131 equations. For a truly final PhD/paper version, the next pass should be a **file-by-file audit of the actual repository at a fixed commit**, recording every script, class, function, dependency, input/output, and exact source location. That would let the report contain genuine `filename.py:function_name`/line-level citations instead of inferred mappings.
